\section{Unfolding and the problem of finding constant series}
\label{sec:unfolding}

\subsection{What a volume computation determines}

Let $\Gamma$ be the mapping class group and let $\gamma$ be a fixed curve
or multicurve class. For an appropriate nonnegative or absolutely integrable
function, the basic unfolding identity is
\begin{equation}
 \int_{\Gamma\backslash\T}
 \sum_{\delta\in\Gamma\gamma}F(\ell_\delta(X))\,\dV(X)
 =\int_{\Stab(\gamma)\backslash\T}
             F(\ell_\gamma(X))\,\dV(X).
 \label{eq:unfold}
\end{equation}
Multiplicity conventions belong in the stabilizer. This identity is just
integration over coset translates of a fundamental domain; it does not
require the curves in the orbit to be simple.

Simplicity becomes important when evaluating the right side. If $\gamma$
is a disjoint simple multicurve, cut along it and use Fenchel--Nielsen
coordinates. Each twist period integrates to $\ell_j$, so the resulting
integral has factors $\prod_j\ell_j\,\dd\ell_j$ times WP volumes of the
components, with the relevant symmetry factors. This is the geometric
integration mechanism in Mirzakhani's volume recursion
\cite[Theorem 8.1]{MirzVolumes}.

Suppose a series $S(X)$ is independently known to be constant, say $C$.
Then
\begin{equation}
 C=\frac{1}{V_{g,n}}
       \int_{\M_{g,n}}S(X)\,\dV(X),
 \label{eq:recoverconstant}
\end{equation}
and unfolding can calculate its value. If constancy has not been proved,
the same quotient computes only its average. This is the precise logical
distinction behind the proposal to recover an identity from a volume.

On the punctured torus, for an unoriented essential simple curve orbit,
the stack-normalized formula is
\begin{equation}
 \int_{\M_{1,1}}\sum_{\alpha\text{ simple}}f(\ell_\alpha)\,\dV
       =\frac12\int_0^\infty x f(x)\,\dd x.
 \label{eq:torusunfold}
\end{equation}
For $f(x)=(1+e^x)^{-1}$, expand in $e^{-x}$ to obtain
\[
 \int_0^\infty\frac{x\,\dd x}{1+e^x}
   =\sum_{m\ge1}\frac{(-1)^{m-1}}{m^2}=\frac{\pi^2}{12}.
\]
The integral of the series is $\pi^2/24$; division by
$V_{1,1}=\pi^2/12$ gives $C=1/2$. Its pointwise constancy is McShane's
theorem~\cite{McShane}:
\begin{equation}
 \sum_{\alpha\text{ simple, unoriented}}
                 \frac{1}{1+e^{\ell_\alpha(X)}}=\frac12.
 \label{eq:mcshane}
\end{equation}
The tetrahedral volume calculation and the covering argument provide an
independent route to the volume used in this normalization check.

\subsection{Operations that preserve constancy}

There is indeed a large closure of any known collection of identities.
Finite linear combinations preserve constancy. So do absolutely convergent
parameter integrals and locally uniformly convergent differentiations of
identities on an identified parameter family. Products of identities give
constant double series when Tonelli's theorem or absolute convergence
applies. These are useful mechanisms for constructing new kernels.

Their geometric interpretation needs to be retained. A product usually
indexes pairs of curves that may intersect; it is not automatically a
disjoint-multicurve sum. A parameter derivative also differentiates lengths
if the metric varies with that parameter. In particular, differentiating
a constant series with respect to the metric gives the one-form identity
\begin{equation}
 \sum_\gamma F'(\ell_\gamma)\,\dd\ell_\gamma=0,
 \label{eq:derivativeidentity}
\end{equation}
assuming sufficient local uniform control. It does not generally give
another scalar identity by differentiating only the displayed weight.

Two possible tests for a proposed constant series follow. One can try to
show that~\eqref{eq:derivativeidentity} vanishes on a spanning set of tangent
directions. Alternatively, if $S\in L^2$ and its mean is $\overline S$, then
\begin{equation}
 \int_{\M}(S-\overline S)^2\,\dV=0
 \label{eq:variance}
\end{equation}
proves constancy almost everywhere, and everywhere if $S$ is continuous.
The second-moment calculation involves pairs of curves, including their
intersection patterns. It can be substantially harder than the first
unfolding integral.

Thus a useful search for further identities should ask for more than
formal closure: a natural index set, a manageable positive weight, a
geometric partition or telescoping mechanism, and a convergence proof.
The constructions in Sections~\ref{sec:figures} and~\ref{sec:exacttwist}
illustrate these requirements.

\subsection{Other established constant series}

There are also constant series with substantially different kernels.
For example, write
\[
 \mathcal L(x)=\operatorname{Li}_2(x)+\tfrac12\log x\log(1-x)
\]
for the Rogers dilogarithm on $(0,1)$, extended continuously to its
endpoints, and define
\begin{align*}
 D(r)={}&\mathcal L\bigl(\sech^2(r/2)\bigr)\\
 &+2\mathcal L\left(\frac{1+e^{-r}}2\right)
  -2\mathcal L\left(\frac{1-e^{-r}}2\right).
\end{align*}
Hu--Tan's punctured-torus identity is
\begin{equation}
 \sum_{\alpha\text{ simple, unoriented}}D(\ell_\alpha)=\frac{\pi^2}{2}.
 \label{eq:dilogtorus}
\end{equation}
The same paper treats a positive torus boundary~\cite[Theorems 1--2]{HuTan}.

For a closed surface of genus $g\ge2$, Luo--Tan give explicit dilogarithmic
weights $f(P),g(T)$ on embedded geometric pairs of pants and one-holed tori
such that
\begin{equation}
 \sum_P f(P)+\sum_T g(T)=8\pi^2(g-1)=\operatorname{Vol}(T^1X).
 \label{eq:luotan}
\end{equation}
Their construction partitions unit tangent directions~\cite{LuoTan}.
It illustrates another source of constant series: the invariant measure
being partitioned need not live only on a horocycle or a geodesic boundary.
The precise, longer subsurface weights are given in that reference.

